\documentclass{article}
\usepackage[utf8]{inputenc}
\usepackage{amsmath, amssymb}
\usepackage{siunitx}
\usepackage{booktabs}
\title{Detaljerte beregninger for EMP-generator (revidert)}
\author{}
\date{}

\begin{document}
\maketitle

\section{Innledning}
Dette dokumentet inneholder alle beregninger som ligger til grunn for designet. Hver beregning er forklart steg for steg, med enheter og konstanter. Vi skiller mellom teoretiske maksimumsverdier og realistiske estimater med dokumenterte tap.

\section{Hovedkondensator og spole}

\subsection{Spesifikasjoner}
\begin{itemize}
\item Kondensator: $C = 1000\,\mu\text{F}$, merkespenning $450\,\text{V}$ (lades til $V_0 = 400\,\text{V}$ for sikkerhetsmargin).
\item Spole: Flat spiral, $N = 6$ vindinger, indre diameter $20\,\text{mm}$, ytre diameter $100\,\text{mm}$, tråd: $10\,\text{AWG}$ ($2.588\,\text{mm}$ diameter).
\end{itemize}

\subsection{Induktans for flat spiralspole (korrigert)}
Wheelers formel for en flat spiralspole i luft, med dimensjoner i tommer og resultat i µH, er:
\begin{equation}
L = \frac{r^2 N^2}{8r + 11w},
\end{equation}
der $r$ er middelradius i tommer, $w$ er viklingens bredde i tommer.

Konvertering: $r = 30\,\text{mm} = 1.181\,\text{tommer}$, $w = 40\,\text{mm} = 1.575\,\text{tommer}$.
\begin{equation}
L = \frac{(1.181)^2 \cdot 6^2}{8 \cdot 1.181 + 11 \cdot 1.575} = \frac{1.395 \cdot 36}{9.448 + 17.325} = \frac{50.22}{26.773} \approx 1.876\,\mu\text{H}.
\end{equation}
Vi bruker $L = 1.9\,\mu\text{H}$ i de videre beregningene.

\subsection{Energi lagret i kondensator}
\begin{equation}
E = \frac{1}{2} C V_0^2 = 0.5 \cdot 1000 \times 10^{-6} \cdot (400)^2 = 80\,\text{J}.
\end{equation}

\section{Toppstrøm og pulsform}

\subsection{Karakteristisk impedans}
\begin{equation}
Z_0 = \sqrt{\frac{L}{C}} = \sqrt{\frac{1.9 \times 10^{-6}}{1000 \times 10^{-6}}} = \sqrt{1.9 \times 10^{-3}} \approx 0.0436\,\Omega.
\end{equation}

\subsection{Toppstrøm (ideell, tapsfri)}
\begin{equation}
I_{\text{peak, ideal}} = \frac{V_0}{Z_0} = \frac{400}{0.0436} \approx 9\,174\,\text{A}.
\end{equation}

\subsection{Demping og reell toppstrøm}
Med total motstand $R \approx 80\,\text{m}\Omega$ (ESR + lysbue + ledninger) blir dempingsfaktoren:
\begin{equation}
\alpha = \frac{R}{2L} = \frac{0.08}{2 \cdot 1.9 \times 10^{-6}} \approx 21\,050\,\text{s}^{-1}.
\end{equation}
Resonansvinkelfrekvens:
\begin{equation}
\omega_0 = \frac{1}{\sqrt{LC}} = \frac{1}{\sqrt{1.9 \times 10^{-6} \cdot 1000 \times 10^{-6}}} = \frac{1}{\sqrt{1.9 \times 10^{-9}}} \approx 22\,940\,\text{rad/s}.
\end{equation}
Frekvens: $f = \omega_0/(2\pi) \approx 3.65\,\text{kHz}$.
Siden $\alpha \ll \omega_0$, er dempingen moderat. Toppstrømmen reduseres til ca. $8\,000$--$9\,000$ A i praksis.

\subsection{Pulsvarighet}
Periode: $T = 1/f \approx 274\,\mu\text{s}$. Tidskonstant for innhyllingskurven: $\tau = 1/\alpha \approx 47.5\,\mu\text{s}$. Etter $5\tau \approx 240\,\mu\text{s}$ er svingningene nesten borte.

\section{Magnetfelt}

\subsection{Felt i sentrum av spolen}
\begin{equation}
B = \frac{\mu_0 N I}{2 r}.
\end{equation}
Med $N=6$, $r=0.05\,\text{m}$, $I=9\,000\,\text{A}$:
\begin{equation}
B = \frac{4\pi \times 10^{-7} \cdot 6 \cdot 9000}{2 \cdot 0.05} \approx 0.679\,\text{T}.
\end{equation}

\subsection{Feltets avtagning med avstand}
For $d = 10\,\text{cm}$:
\begin{equation}
B(0.1) \approx 0.679 \left( \frac{0.05}{\sqrt{0.05^2 + 0.1^2}} \right)^3 \approx 0.679 \cdot 0.089 \approx 0.060\,\text{T}.
\end{equation}

\section{Indusert spenning i test-sløyfe}
Med $N_p=5$, $A_p=3.14\times10^{-4}\,\text{m}^2$, $\Delta t \approx T/4 = 68.5\,\mu\text{s}$:
\begin{equation}
\mathcal{E} \approx -5 \cdot 3.14\times10^{-4} \cdot \frac{0.060}{68.5\times10^{-6}} \approx -1.38\,\text{V}.
\end{equation}
I virkeligheten kan $dB/dt$ være høyere i det øyeblikket strømmen stiger raskest, så spenningen kan bli 5–10 V.

\section{Termiske beregninger}

\subsection{Oppvarming av spole}
Energi per puls: $80\,\text{J}$. Massen av spolen: $N=6$ vindinger, hver med lengde $\approx \pi \cdot 0.1\,\text{m} = 0.314\,\text{m}$, total lengde $1.884\,\text{m}$. Tråd: $10\,\text{AWG}$, diameter $2.588\,\text{mm}$, tverrsnitt $A_{\text{tråd}} = \pi (1.294\times10^{-3})^2 = 5.26 \times 10^{-6}\,\text{m}^2$. Volum: $V = 1.884 \cdot 5.26\times10^{-6} = 9.91 \times 10^{-6}\,\text{m}^3$. Tetthet kobber: $8960\,\text{kg/m}^3$, masse $m = 8960 \cdot 9.91\times10^{-6} = 0.0888\,\text{kg} = 88.8\,\text{g}$. Varmekapasitet kobber: $c = 385\,\text{J/(kg·K)}$. Temperaturøkning per skudd:
\begin{equation}
\Delta T = \frac{E}{m c} = \frac{80}{0.0888 \cdot 385} \approx \frac{80}{34.2} \approx 2.34\,\text{K}.
\end{equation}
Dette er helt ufarlig for spolen. Ved gjentatte skudd må man likevel passe på at varmen ledes bort.

\subsection{Oppvarming av kondensator}
Kondensatoren har en intern motstand (ESR) på typisk $50\,\text{m}\Omega$. Under utladning vil en del av energien gå tapt som varme i kondensatoren. Tapet kan estimeres ved å anta at strømmen deles mellom ESR og lasten. Siden lastmotstanden (spole + gnistgap) er i samme størrelsesorden, kan omtrent halvparten av energien tapes i kondensatoren, dvs. $40\,\text{J}$. En 1000 µF kondensator veier ca. $50\,\text{g}$ og har en varmekapasitet på ca. $50\,\text{J/K}$ (aluminium). Temperaturøkning: $\Delta T \approx 40/50 = 0.8\,\text{K}$ per skudd. Også akseptabelt.

\section{Komponentpåkjenninger}

\subsection{MOSFET i ladekrets}
Under avslag kan drain-spenningen nå opp til $2 \times V_{\text{bat}} + V_{\text{reflektert}}$. Med $V_{\text{bat}} = 54\,\text{V}$ og et viklingsforhold $n = 200/10 = 20$, er reflektert spenning $V_{\text{ut}}/n = 400/20 = 20\,\text{V}$. Total spenning: $2 \cdot 54 + 20 = 128\,\text{V}$. IRF840 tåler $500\,\text{V}$, så marginen er god.

\subsection{Dioder}
UF4007 tåler $1000\,\text{V}$ repeterende. I verste fall kan utgangsspenningen nå $450\,\text{V}$, så marginen er over 2x.

\subsection{Gnistgap}
Wolframelektroder har et smeltepunkt på $3422\,^{\circ}\text{C}$ og tåler ekstremt høye pulsstrømmer uten å erodere nevneverdig.

\section{Konklusjon}
Designet er robust, men stigetiden er relativt lang (~70 µs). For raskere puls og høyere induserte spenninger kan en peaking-kondensator vurderes (se uløste problemer).
\end{document}
